\section{Parameter scaling of the constructed ideals}

\begin{lemma}[Actual parameter substitutions]
\label{lem:pl-parameter-scaling-hom}
\uses{lem:pl-fs-product-ideal}
\lean{EnvelopeProperty.parameterScaleHom,EnvelopeProperty.parameterScaleHom_C,EnvelopeProperty.parameterScaleHom_X,EnvelopeProperty.parameterScaleHom_comp,EnvelopeProperty.parameterScaleHom_one,EnvelopeProperty.parameterMonomial,EnvelopeProperty.parameterScale_map_base,EnvelopeProperty.parameterScale_monomial}
\leanok
For every commutative coefficient ring $R$ and every $c\in R$, substitution
$t\mapsto ct$ defines an actual $R$-algebra endomorphism of $R[t]$.
These substitutions compose by multiplication and the substitution for
$c=1$ is the identity.  Acting on the coefficients of
$R[t][x_0,\ldots,x_M]$, it sends $s(x)t^e$ to $c^e s(x)t^e$.
\end{lemma}
\begin{proof}
The actual polynomial evaluation homomorphism sending $t$ to $ct$ and
fixing $R$ gives the substitution.  Two such algebra homomorphisms are
equal when they agree on the coefficients and on $t$; this proves the
composition and identity laws.  Coefficient mapping fixes the polynomial
$s(x)$, while $(ct)^e=c^et^e$, which gives the monomial identity.
\end{proof}

\begin{lemma}[Invariant vertical ideals of finite FS functions]
\label{lem:pl-finite-fs-invariant-vertical-ideal}
\uses{lem:pl-fs-invariant-product-ideal,lem:pl-vertical-truncation,lem:pl-parameter-monomial-ideal-invariance}
\lean{EnvelopeProperty.ProjSubvariety.IsFiniteFS.exists_parameterInvariant_vertical_productIdeal_model}
\leanok
A rational finite FS function constructs a genuine vertical homogeneous
product ideal with the exact integral model-logarithm formula, an
explicit positive parameter power, and a retained finite homogeneous
parameter-monomial generating family.  Its entire ideal is invariant
after every coefficient extension and every unit parameter scaling.
\end{lemma}
\begin{proof}
Use the retained family from the preceding FS construction.  Compactness
of the variety gives $A>0$ with $-A$ below the model logarithm everywhere.
Adjoin the actual parameter monomial $t^A$ to the retained family.
The new ideal logarithm is the maximum of the original logarithm and
$-A$, hence is unchanged on the variety.  The added base form is $1$,
homogeneous of degree zero, so the enlarged finite family remains a
homogeneous parameter-monomial family.  Apply universal coefficient
extension and unit-scaling invariance to this very family.
\end{proof}

\begin{lemma}[Invariant vertical ideals for actual PL functions]
\label{lem:pl-invariant-vertical-ideals}
\uses{lem:pl-finite-fs-invariant-vertical-ideal,prop:pl-vertical-ideal-models}
\lean{EnvelopeProperty.ProjSubvariety.IsPL.exists_parameterInvariant_vertical_productIdeal_models}
\leanok
Every actual PL function constructs two genuine finitely generated vertical
homogeneous product ideals.  Both entire ideals are invariant under all
parameter-unit scalings after every coefficient extension.  Their finite
intrinsic Gauss logarithms recover a positive integral multiple of the
original PL function, up to an integral constant.
\end{lemma}
\begin{proof}
Write the PL function as a difference of two actual rational finite FS
functions.  The preceding construction yields ideals $I_0,J_0$, positive
integers $r_1,r_2$, and constants $C_1,C_2$ with
$\log I_0=r_1\alpha-C_1$ and $\log J_0=r_2\beta-C_2$.
Use the actual powers $I=I_0^{r_2}$ and $J=J_0^{r_1}$.
Ideal image commutes with positive integral powers, so the universal
coefficient-extended invariance of each original ideal proves the same
whole-ideal equality for its power.  The original positive parameter powers
give actual positive parameter powers in $I,J$.  Intrinsic ideal logarithms
multiply by the exponent, and both remain finite.  Subtracting gives
$r_1r_2(\alpha-\beta)=\log I-\log J+r_2C_1-r_1C_2$.
\end{proof}

\begin{lemma}[Actual invariant ideals constructed from FS data]
\label{lem:pl-fs-invariant-product-ideal}
\uses{lem:pl-fs-product-ideal,lem:pl-parameter-monomial-ideal-invariance}
\lean{EnvelopeProperty.ProjBerk.FSData.exists_parameterInvariant_productIdeal_model}
\leanok
Rational FS data construct an actual finitely generated homogeneous
product ideal whose whole intrinsic Gauss logarithm is the prescribed
positive integral multiple of the FS function, up to an integral constant.
The ideal is invariant under every coefficient-unit scaling of the
parameter after every coefficient extension, and has an explicitly
retained finite parameter-monomial generating family with homogeneous
base forms.
\end{lemma}
\begin{proof}
Clear the rational weights by $N>0$, choose $C\ge N\lambda_j$, and put
$e_j=C-N\lambda_j$.  Retain the actual finite family
$(s_j^N,e_j)$ and form the genuine homogeneous ideal generated by its
parameter monomials.  All base forms have degree $N\deg(s_j)$.
The generator Gauss logarithms are
$N(\log|s_j|+\lambda_j)-C$; taking their finite maximum gives the exact
FS formula.  Apply the monomial-generator invariance theorem
to this very retained family, which proves invariance of the entire
constructed ideal after every coefficient extension and every unit
scaling over the extended ring.
\end{proof}

\begin{lemma}[Invariance of the actual monomial-generated ideals]
\label{lem:pl-parameter-monomial-ideal-invariance}
\uses{lem:pl-parameter-scaling-hom}
\lean{EnvelopeProperty.parameterScale_span_monomials_le,EnvelopeProperty.parameterScale_span_monomials_eq,EnvelopeProperty.map_parameterMonomial,EnvelopeProperty.parameterScale_span_baseChanged_monomials_eq}
\leanok
Let $I$ be the actual ideal generated by a family of elements $s(x)t^e$.
For every unit $u\in R^\times$, the whole coefficient-substitution image
of $I$ under $t\mapsto ut$ is exactly $I$.  After every actual coefficient
extension $R\to R'$, the entire extended ideal has the same invariance
under every unit of $R'$.
\end{lemma}
\begin{proof}
Each image generator is its original generator times the coefficient
$u^e$, so the entire image ideal is contained in $I$.  Apply the same
inclusion with $u^{-1}$, and then map that inclusion by $t\mapsto ut$.
The two actual ring maps compose to the identity; thus $I$ is contained
in its image, giving equality.  The statement applies to the full
original ideals and does not follow merely from their Gauss functions
or from central support.  Coefficient extension sends $s(x)t^e$ to the
same parameter monomial with the coefficient-extended base form.
Consequently the extended ideal is again generated by actual parameter
monomials, so the same proof applies over $R'$.
\end{proof}
